\documentclass[11pt, a4paper]{article}

% Pacchetti per la codifica e la lingua
\usepackage[utf8]{inputenc}
\usepackage[T1]{fontenc}
\usepackage[italian]{babel}

% Pacchetti per layout e margini
\usepackage{geometry}
\geometry{a4paper, top=2.5cm, bottom=2.5cm, left=2.5cm, right=2.5cm}

% Pacchetti per tabelle migliorate e grafica
\usepackage{tabularx}
\usepackage{booktabs}
\usepackage{graphicx}
\usepackage{hyperref}

% Dati del documento
\title{\textbf{Relazione Accessibilità Prototipo HCI}}
\author{
    \textbf{Gruppo: DJ} \\
    Ughetto Luca (1145888) \\
    Rey Mattia (1143424) \\
}
\date{\today}

\begin{document}

    \maketitle
    \newpage

    \section{Errori riscontrati con gli strumenti}
    
    \textit{Lorem ipusm}\dots

    \newpage
    \section{Aggiunte per garantire l'accessibilità}
    A lezione abbiamo visto diversi tag e diverse pratiche che se implementate permettono agli utenti con disabilità di accedere al sito in maniera più semplice, fluida e veloce.
    \paragraph{Skip link}
    Abbiamo inserito uno \texttt{skip-link} in ogni pagina. In modo che chi accede alla pagina con un lettore riesca a saltare l'introduzione fino al contenuto principale.
    \paragraph{Zoom}
    Ci siamo assicurati che la pagina risultasse coerente e leggibile fino al livello di zoom $200\%$.
    \paragraph{Coerenza dei tag}
    Non sempre nel design originale usavamo i tag adeguati per le funzioni che coprivano. Ad esempio la barra di ricerca veniva definita con un \texttt{div} invece che con un tag più specifico come un \texttt{form} con un ruolo esplicito. Anche per i \texttt{breadcrumbs} è stata fatta un'operazione simile.

    gerarchia degli h.
    scritte in inglese lette in inglese
    emoji isolate
    aria tag
    labels
    in generale descrizioni accurate dei vari tag
\end{document}